%-------------------------
% Engineering Resumes layout
% Based off of: https://www.reddit.com/r/EngineeringResumes/wiki/
%------------------------

% Airier single-column layout. Implements the same command interface as
% classic.tex (references/INTERFACE.md), so one content file renders under
% either. Input this from a driver, after the driver has set its identity
% macros and before \begin{document}.

\usepackage[
  letterpaper,
  top=0.5in,
  bottom=0.5in,
  left=0.5in,
  right=0.5in
]{geometry}
\usepackage[T1]{fontenc}
\usepackage[utf8]{inputenc}

% Font. The driver may \input a fonts/ file before this one; if it did not,
% fall back to the default. Replaces an earlier \ifResumeTimes boolean, which
% could only ever express two choices.
\providecommand{\ResumeFontName}{}
\makeatletter
\ifx\ResumeFontName\@empty
  % Source Sans Pro -- humanist sans, the default.
\newcommand{\ResumeFontName}{sourcesans}
\usepackage[default]{sourcesanspro}

\fi
\makeatother

\usepackage{enumitem}
\usepackage[hidelinks]{hyperref}
\usepackage{titlesec}

\raggedright
\pagestyle{empty}
\input{glyphtounicode}
\pdfgentounicode=1

\titleformat{\section}
  {\bfseries\large}
  {}
  {0pt}
  {}
  [\vspace{1pt}\titlerule\vspace{-6.5pt}]

\renewcommand\labelitemi{$\vcenter{\hbox{\small$\bullet$}}$}
\setlist[itemize]{itemsep=-2pt, leftmargin=12pt, topsep=6pt, before=\small}

%-------------------------
% Command interface

% Entries are plain paragraphs here, not list items, so the outer subheading
% list is a no-op.
\newcommand{\resumeSubHeadingListStart}{}
\newcommand{\resumeSubHeadingListEnd}{}

\newcommand{\resumeItemListStart}{\begin{itemize}}
\newcommand{\resumeItemListEnd}{\end{itemize}}
\newcommand{\resumeItemListStartTight}{\begin{itemize}}

\newcommand{\resumeItem}[1]{\item #1}
\newcommand{\resumeSubItem}[1]{\item #1}

\newcommand{\resumeSubheading}[4]{%
  \textbf{#1} -- \textit{#3} \hfill #2 \\%
  \vspace{-9pt}%
}

% A degree is not a job title, so it stays upright here even though the job
% title above is italicised.
\newcommand{\resumeEducation}[3]{%
  \textbf{#1} -- #3 \hfill #2 \\%
}

% Publication titles are long, so the status note gets its own line rather
% than being appended to the title the way a job title is.
\newcommand{\resumePublication}[3]{%
  \textbf{#1} \hfill #2 \\%
  \textit{#3} \\%
  \vspace{-9pt}%
}

\newcommand{\resumeSubSubheading}[2]{%
  \textit{#1} \hfill #2 \\%
  \vspace{-9pt}%
}

\newcommand{\resumeProjectHeading}[2]{%
  \textbf{#1} \hfill #2 \\%
  \vspace{-9pt}%
}

\newcommand{\resumeProjectTitle}[1]{\item[] \textbf{#1}}

\newcommand{\resumeSectionGap}{\vspace{-18.5pt}}

% Negative leading tuned for a denser layout does not transfer to these
% metrics -- honouring it overlaps headings with body text. Ignore it.
\newcommand{\resumeTighten}[1]{}

% Identity defaults. Drivers set these via profiles/<person>.tex.
\providecommand{\ResumeName}{Your Name}
\providecommand{\ResumeEmail}{you@example.com}
\providecommand{\ResumePhone}{000-000-0000}
\providecommand{\ResumeLinkedInName}{Your Name}
\providecommand{\ResumeLinkedInURL}{https://www.linkedin.com/in/your-handle/}

%----------HEADING----------
\newcommand{\ResumeHeader}{%
\centerline{\Huge \ResumeName}
\vspace{5pt}
\centerline{
  \href{mailto:\ResumeEmail}{\ResumeEmail}
  \enspace|\enspace \ResumePhone
  \enspace|\enspace
  \href{\ResumeLinkedInURL}{LinkedIn: \ResumeLinkedInName}
}
\vspace{-10pt}
}
